
\slajdid{05-s019}
\begin{frame}[label=05-s019]
\gusttekst\setlength{\parskip}{1pt}\setlength{\abovedisplayskip}{4pt}\setlength{\belowdisplayskip}{4pt}
Pretpostavke:\\
1. radimo sa fiksnom tačkom i ignorišemo njenu poziciju,\\
\hspace*{5mm}smatramo da su operandi usaglašeni sa pozicijama tačke pa radimo sa celobrojnim vrednostima\\
2. radimo sa ograničenim brojem cifara $n$ i za operande i za rezultate\\
3. imamo u registru potrebne vodeće i prateće nule
\[a_i,b_i,s_i\in\{0,1,2,\ldots r-2,r-1\}\]
\begin{columns}[c,onlytextwidth]
\column{.38\textwidth}\centering\begin{tikzpicture}[x=.70cm,y=.48cm,font=\fontsize{9}{10}\selectfont]
\foreach \x/\i in {0/{n-1},1/{n-2},3/i,5/1,6/0}{
\node at (\x,1) {$a_{\i}$};\node at (\x,0) {$b_{\i}$};\node at (\x,-1) {$s_{\i}$};}
\foreach \x in {2,4}{\foreach \y in {-1,0,1}{\node at (\x,\y) {$\ldots$};}}
\node[DEcrvena] at (-.7,-1) {$s_n$};\draw[DEvod] (-.95,-.5)--(6.4,-.5);
\foreach \x in {0,1,2,3.5,5}{\draw[DEplava,-{Latex[length=1.4mm]}] (\x-.15,2)--(\x-.4,2)--(\x-.4,1.35);}
\end{tikzpicture}

\column{.59\textwidth}
B – Borrow, pozajmica 0 ili 1 bez obzira na osnovu brojnog sistema

Oduzimanje radimo u osnovi brojnog sistema
\end{columns}
\[(a_i-B_{i-1}-b_i)_{min}=0-1-(r-1)=\underset{\substack{\textcolor{DEplava}{\uparrow}\\\displaystyle B_i=1}}{-r}+\underset{\substack{\textcolor{DEplava}{\uparrow}\\\displaystyle s_i=0}}{0}\]
\begin{columns}[c,onlytextwidth]
\column{.57\textwidth}
\[\left.\begin{gathered}s_n=0-0-B_{n-1}=0\ \mathit{ili}\ -1\quad\textcolor{DEcrvena}{???}\\[2mm]
s_n=-1\quad\text{\textcolor{DEcrvena}{\makecell{Prekoračenje opsega, rezultat nije\\moguće smestiti u $n$ cifara???}}}\end{gathered}\right\}\]
\column{.40\textwidth}\centering
Javlja se kada je $a<b$\\odnosno kada rezultat treba da bude negativan\\O ovome kasnije
\end{columns}
\end{frame}
